% Part: first-order-logic
% Chapter: introduction
% Section: models-theories

\documentclass[../../../include/open-logic-section]{subfiles}

\begin{document}

\olfileid{fol}{int}{mod}

% SEGMENT OLP-0147-S01
\olsection{மாதிரிகளும் கோட்பாடுகளும்}

முதல்தர தருக்கத்தின் தொடரமைப்பையும் பொருண்மையியலையும் வரையறுத்துவிட்டதும்,
!!{structure}s மற்றும் பொருண்மையியல் கருத்துகளின் பண்புகளை ஆராயும் பணியைத்
தொடங்கலாம். !!{derivation} முறைமைகளையும் வரையறுத்து அவற்றையும் ஆராயலாம்.
!!{sentence}s கொண்ட ஒரு கணத்திற்கு, அந்தக் கணத்திலுள்ள எல்லா !!{sentence}s ஐயும்
மெய்யாக்கும் !!{structure}s எவை என்று கேட்கலாம். !!{sentence}s கணம்~$\Gamma$
ஒன்று கொடுக்கப்பட்டால், அவற்றை நிறைவுறுத்தும் !!a{structure}~$\Struct{M}$
``$\Gamma$ இன் \emph{மாதிரி}'' எனப்படுகிறது. $\Gamma$ இலிருந்து தொடங்கி
அதன் மாதிரிகளைத் தேடலாம்---அவை எவ்வாறு இருக்கும்? அவை எவ்வளவு பெரியதாகவோ
சிறியதாகவோ இருக்க வேண்டும்? ஆனால் ஒரு !!{structure} அல்லது !!{structure}s
தொகுப்பிலிருந்து தொடங்கி, அவற்றில் எந்த !!{sentence}s மெய்யாக இருக்கின்றன
என்றும் கேட்கலாம். இவ்வ!!{structure}s ஐ, அவற்றில் மட்டுமே அவை மெய்யாக இருக்கும்
வகையில், ஏதேனும் !!{sentence}s \emph{வரையறுக்கின்றனவா}? இவ்வகைக் கேள்விகள்
\emph{மாதிரிக் கோட்பாட்டின்} களமாகும். அவை \emph{அடிகோள் முறை}க்கும்
அடித்தளமாக உள்ளன: ஒரு கோட்பாட்டின் அடிகோள்களான !!{sentence}s கணத்தால்
!!{structure}s தொகுப்பை விவரித்தல். அடிகோள்களால் முதல்தர தருக்கத்தில்
உட்கிடையாகும் !!{sentence}s மட்டுமே அடிகோள்களின் எல்லா மாதிரிகளிலும்
மெய்யாகின்றன என்ற கவனிப்பே இதைச் சாத்தியமாக்குகிறது.

% SEGMENT OLP-0147-S02
மிக எளிய எடுத்துக்காட்டாக, முன்வரிசைகளைக் கருதுக. முன்வரிசை என்பது ஏதோ ஒரு
கணம்~$A$ மீது உள்ள தொடர்பு~$R$; அது தற்சுட்டானதாகவும் கடப்புநிலையானதாகவும்
இருக்கும். $A$ மீது $R \subseteq A \times A$ என்ற இரு-இடத் தொடர்பு கொண்ட
கணம், ஒரே இரு-இடத் தொடர்புக் குறி~$\Obj P$ கொண்ட முதல்தர மொழிக்கான
!!a{structure} ஒன்றை வழங்கத் தேவையானதே: $\Domain{M} = A$ என்றும்
$\Assign{\Obj P}{M} = R$ என்றும் வைப்போம். $R$ முன்வரிசை என்பதால் அது
தற்சுட்டானதும் கடப்புநிலையானதுமாகும்; இதைக் கூறும் முதல்தர தருக்கத்தின்
!!{sentence}s கணம்~$\Gamma$ ஒன்றைக் கண்டறியலாம்:
\begin{align*}
  & \lforall[\Obj v_0][\Atom{\Obj P}{\Obj v_0,\Obj v_0}]\\
  & \lforall[\Obj v_0][\lforall[\Obj v_1][\lforall[\Obj v_2][((\Atom{\Obj P}{\Obj v_0,\Obj v_1} \land \Atom{\Obj P}{\Obj v_1,\Obj v_2}) \lif \Atom{\Obj P}{\Obj v_0,\Obj v_2})]]]
\end{align*}
இவ்வ!!{sentence}s, ``எந்த~$x$ க்கும் $Rxx$'' ($R$ தற்சுட்டானது) மற்றும்
``$Rxy$ மற்றும் $Ryz$ எப்போதெல்லாம் உள்ளனவோ அப்போதெல்லாம் $Rxz$ உம்''
($R$ கடப்புநிலையானது) என்பவற்றின் குறியாக்கங்கள் மட்டுமே. இந்த இரு
!!{sentence}s கொண்ட~$\Gamma$ இன் மாதிரி !!a{structure}~$\Struct{M}$ என்பது அப்போதும் அப்போது
மட்டுமே $R$ (அதாவது, $\Assign{\Obj P}{M}$) ஆனது~$A$ மீது ஒரு முன்வரிசை
($\Domain{M}$) ஆகும். வேறு சொற்களில்,~$\Gamma$ இன் மாதிரிகள் சரியாக
முன்வரிசைகளே. $\Obj P$ ஐ !!{predicate} ஆகக் கொண்ட முதல்தர மொழியில்
வெளிப்படுத்தக்கூடிய எல்லா முன்வரிசைப் பண்புகளும் (மேலுள்ள தற்சுட்டுத் தன்மை,
கடப்புநிலை போன்றவை)~$\Gamma$ இன் இரு !!{sentence}s களால் உட்கிடையாகின்றன;
மாறாகவும் அதே. எனவே~$\Gamma$ இன் மாதிரிகளைப் பற்றி நாம் நிறுவக்கூடிய எதையும்
எல்லா முன்வரிசைகளைப் பற்றியும் நிறுவியுள்ளோம்.

% SEGMENT OLP-0147-S03
எந்தக் குறிப்பிட்ட கோட்பாட்டுக்கும் மாதிரிகளின் வகுப்புக்கும் (உதாரணமாக,
$\Gamma$ மற்றும் எல்லா முன்வரிசைகள்), அதற்குரிய முதல்தர மொழியில் எதை
வெளிப்படுத்த முடியும், எதை முடியாது என்பவை பற்றிய ஆர்வமூட்டும் கேள்விகள்
இருக்கும். எல்லா மொழிகளுக்கும் மாதிரி வகுப்புகளுக்கும் சுவாரசியமான சில
!!{structure} பண்புகள், அவற்றின் !!{domain} அளவைப் பற்றியவை. உதாரணமாக,
எந்த~$n \in \PosInt$ க்கும் !!{domain} இல் சரியாக~$n$ !!{element}s உள்ளன
என்று வெளிப்படுத்தலாம். முடிவிலி எண்ணிக்கையிலான !!{sentence}s ஐப் பயன்படுத்தி
!!{domain} முடிவிலியானது என்றும் வெளிப்படுத்தலாம். ஆனால் பொருட்களம்
இறுதிக்குட்பட்டது என்றும், அல்லது !!{nonenumerable} என்றும் வெளிப்படுத்த
முடியாது. முதல்தர மொழிகளின் வரம்புகள் பற்றிய இம்முடிவுகள் இறுதிக்குட்பட்ட
தன்மை மற்றும் லோவென்ஹைம்--ஸ்கோலம் தேற்றங்களின் விளைவுகள்.

\end{document}
